Style augmentation is a widely used way to close the domain gap in single-source cross-domain classification, but its reported benefit is confounded: a styled training set differs from the original in its size, in its generic appearance diversity, and in the style it carries, all at once. We measure what style augmentation actually buys on PACS with four controlled experiments. First, two data-volume controls --- a byte-identical duplicate of every training image, and a non-stylistic transform of every training image --- separate the effect of the extra rows from the effect of the style. Second, because a spatial strength map necessarily applies less total style than a uniform mix at the same nominal strength, we train a strength-matched uniform control and an equal-strength spatial variant, so that the effect of the spatial distribution can be isolated. Third, a sweep over the style strength traces the dose--response relation, and the style pool is split into disjoint halves to test the robustness of the result. Adding AdaIN-styled images raises the cross-domain average from 40.8\% to 47.4\% over 5 seeds, a paired gain of 6.5 points (95\% CI [3.1, 9.9], $p=0.0060$); neither control reproduces an appreciable part of it (+Duplicate -3.1 points, $p=0.031$; +NonStyle +2.2 points over +Duplicate, $p=0.047$). Under the strength-matched control the spatial modulation and the uniform application are indistinguishable (-0.17 points, $p=0.775$), and the accuracy follows the average style strength rather than its spatial distribution. We therefore conclude that style augmentation helps because of the style it applies, and that the benefit of spatially adaptive stylization is not supported by these controlled experiments.
